%--------------------------------------------------------------------------------
%
%
%--------------------------------------------------------------------------------
\chapter*{Interruptions}
\addcontentsline{toc}{chapter}{Interruptions}

Twin-64 features a simple but powerful interrupt system. Interruptions are grouped in to four different classes, faults, traps, interrupts, and checks that may happen during instruction processing:

\begin{itemize} 

\item{Fault - The current instruction requests a legitimate action which cannot be carried out due to a system problem, such as the absence of a page from main memory. After the system problem has been corrected, the faulting instruction will execute normally. Faults are synchronous with respect to the instruction stream.}

\item{Trap - Traps include two sorts of possibilities: either the function requested by the current
instruction cannot or should not be carried out, or system intervention is desired by
the user before or after the instruction is executed. Examples of the first type include arithmetic operations that result in signed overflow and instructions executed with insufficient privilege for their intended function. Such instructions are normally not re-executed. Examples of the second type include the debugging support traps. Traps are synchronous with respect to the instruction stream}

\item{Interrupt - An external entity (for example, an I/O device or the power supply) requires attention. Interrupts are asynchronous with respect to the instruction stream.}

\item{ Check - The processor has detected an internal malfunction. Checks can be either
synchronous or asynchronous with respect to the instruction stream.
All four classes of interruptions are handled in the same way.}

\end{itemize}


\section*{Interrupt Vector Table}
\addcontentsline{toc}{section}{Interrupt Vector Table}

The control register IVA holds the virtual address of the interrupt vector table. Upon an interrupt, the current instruction address and the instruction is saved to the respective control registers and control branches then to the interrupt handler in the interrupt vector table. Each entry is just a set of instructions that can be executed.

The interruptions are categorized into four
groups based on their priorities:

\begin{longtable}{@{}p{0.075\linewidth}p{0.25\linewidth}p{0.2\linewidth}p{0.15\linewidth}p{0.15\linewidth}@{}}
    \caption{Trap Table} \\
    \toprule
    \textbf{Trap} & \textbf{Name} & \textbf{Instruction Adr} & \textbf{Arg0} & \textbf{Arg1}\\
    \midrule
    \endfirsthead
    \toprule
    \textbf{Trap} & \textbf{Name} & \textbf{IA} & \textbf{Arg0} & \textbf{Arg1}\\
    \midrule
    \endhead
    \midrule
    \multicolumn{2}{r}{\textit{Continued on next page}} \\
    \midrule
    \endfoot
    \bottomrule
    \endlastfoot
    \textbf{1} & Machine Check & current instruction & check type & - \\
    \midrule
    \textbf{2} & Power Failure & current instruction & - & - \\
    \midrule
    \textbf{3} & Recovery Counter & current instruction & - & - \\
    \midrule
    \textbf{4} & External Interrupt & next instruction & - & - \\
    \midrule
    \textbf{5} & Illegal instruction & current instruction & instruction & - \\
    \midrule
    \textbf{6} & Privileged operation & current instruction & instruction & - \\
    \midrule
    \textbf{7} & Privileged register & current instruction & instruction & - \\
    \midrule
    \textbf{8} & Integer Overflow& current instruction & instruction & - \\
    \midrule
    \textbf{9} & Instruction TLB miss& current instruction & - & - \\
    \midrule
    \textbf{10} & Non-access Instruction TLB miss & current instruction & - & - \\
    \midrule
    \textbf{11} & Instruction memory protection & current instruction & - & - \\
    \midrule
    \textbf{12} & Data TLB miss& current instruction & instruction & target address \\
    \midrule
    \textbf{13} & Non-access Data TLB miss & current instruction & instruction & target address \\
    \midrule
    \textbf{14} & Data memory access rights & current instruction & instruction & target address \\
    \midrule
    \textbf{15} & Data memory protection & current instruction & instruction & target address \\
    \midrule
     \textbf{16} & Page reference & current instruction & instruction & target address \\
    \midrule
    \textbf{17} & Alignment & current instruction & instruction & target address \\
    \midrule
    \textbf{18} & Break Instruction & current instruction & instruction & - \\
    \midrule
     \textbf{19..31} & reserved & - & - & - \\
    \midrule
\end{longtable}

\section*{Interruption Handling}
\addcontentsline{toc}{section}{Interruption Handling}

mask and request

EIR in I/O space ...

EIM in CR

status bit in PSW





